\chapter{Integer algorithms}
\MARSChapterHeading{Integer algorithms}
\label{revision:integer-algorithms}

A line between two pixels rarely passes through the centers of all the pixels
between them. Drawing it therefore requires a decision: which nearby cells should
represent the continuous line? The decision must work for a steep line, a falling
line and a line consisting of one point. It should also give the same image when
the caller exchanges the two endpoints.

We will draw on an eight-by-eight word framebuffer. This is smaller than the
previous chapter's paired sprite display: a row now occupies 32 bytes, and the
whole frame occupies 256 bytes. Existing pixels outside the selected line must
survive. The result is an overlay, not a cleared image followed by a drawing.

You already know how to calculate a pixel address, use signed comparisons and
write a leaf function. This chapter adds an integer decision state. You should be
able to explain its update rule, distinguish a saved decision from the state it
changes, account for endpoint and tie cases, and adapt the raster to a related
operation without changing its geometry.

\section{Specify the cells before writing the loop}

Let the first endpoint be $(1,1)$ and the second $(6,4)$. Coordinates use the screen
convention: $x$ increases to the right and $y$ increases downward. Five horizontal
steps cover three vertical steps. At $x=2$, the continuous line has height $1.6$,
so row two is nearer than row one. At $x=3$, its height is $2.2$, so row two remains
nearer. Continuing gives the six cells in Figure~\ref{fig:revision-line-raster}.
Both endpoints are included; five horizontal steps visit six columns.

\begin{figure}[htbp]
\centering
\begin{tikzpicture}[x=6mm,y=-6mm,alt={An eight by eight grid has x increasing right and y increasing downward. The continuous line joins centers one comma one and six comma four. Its six selected cells are one comma one, two comma two, three comma two, four comma three, five comma three, six comma four.}]
\foreach \x/\y in {1/1,2/2,3/2,4/3,5/3,6/4} {
  \fill[black!18] (\x,\y) rectangle ++(1,1);
}
\foreach \n in {0,...,8} {
  \draw[black!45,thin] (\n,0) -- (\n,8);
  \draw[black!45,thin] (0,\n) -- (8,\n);
}
\draw[black,dashed,thick] (1.5,1.5) -- (6.5,4.5);
\foreach \x/\y in {1/1,2/2,3/2,4/3,5/3,6/4} {
  \fill (\x+0.5,\y+0.5) circle[radius=0.7mm];
}
\foreach \n in {0,...,7} {
  \node[font=\small,anchor=south] at (\n+0.5,-0.2) {\n};
  \node[font=\small,anchor=east] at (-0.2,\n+0.5) {\n};
}
\node[anchor=west] at (8.3,-0.2) {$x$};
\node[anchor=east] at (-0.2,8.4) {$y$};
\end{tikzpicture}
\caption{The dashed continuous line and the selected cell centers serve different purposes. Each shaded cell receives one color word; all other words are preserved.}
\label{fig:revision-line-raster}
\end{figure}


At an exact halfway position, two rows are equally near. We need a convention to
make the result reproducible. Order the endpoints by increasing $x$, breaking an
equal-$x$ tie by increasing $y$. Call the resulting endpoints $A$ and $B$. At a
halfway choice, take the minor-coordinate step toward $B$. Thus $(0,0)$ to $(2,1)$
selects $(1,1)$ between its endpoints. This ordering makes the halfway choice independent of the supplied direction.

For a shallow line, write $d_x=B_x-A_x$ and $D=|B_y-A_y|$, with $d_x\geq D$.
At horizontal displacement $i$, the ideal vertical displacement has magnitude
$iD/d_x$. Rounding to the nearest integer, with halfway values advancing, gives
\[
 j=\left\lfloor\frac{2iD+d_x}{2d_x}\right\rfloor.
\]
Here $\lfloor q\rfloor$ means the greatest integer no larger than $q$. For a steep
line, exchange the roles of the major and minor displacements. These formulas
specify the selected cells. They also provide a way to test the result independently
of the incremental algorithm we will implement. Handle a point separately so that
no formula divides by zero.

\section{Carry the rounding decision forward}

Calculating a fraction for every cell would repeat work. The numerator changes
by a fixed amount each time the cursor advances. An integer accumulator can retain
that change between iterations. This is the central idea of Bresenham line drawing,
originally developed for digital plotters.\footnote{J. E. Bresenham, ``Algorithm for
computer control of a digital plotter,'' \emph{IBM Systems Journal} 4(1), 1965,
pp.~25--30. \url{https://doi.org/10.1147/sj.41.0025}. The implementation here uses
an all-slope recurrence and an explicit endpoint/tie contract.}

Use $d_y=-D$, so $d_y\leq0$, and let $s_y$ be $+1$ when $B_y\geq A_y$, otherwise
$-1$. The horizontal direction is always positive after ordering; a vertical line
will simply take no horizontal step. Initialize the accumulator to
$e=d_x+d_y$. At the currently selected cell, save $e_2=2e$. Then make two separate
decisions:
\[
\begin{aligned}
 e_2\geq d_y &: \quad e\gets e+d_y,\qquad x\gets x+1,\\
 e_2\leq d_x &: \quad e\gets e+d_x,\qquad y\gets y+s_y.
\end{aligned}
\]
Either coordinate can advance, and both can advance in the same iteration. Paint
the current cell and test endpoint equality before applying those updates.

To see where the second comparison comes from, consider the shallow case. After
$i$ horizontal and $j$ vertical steps, define $E=iD-jd_x$. The next ideal height
requires a vertical step when $(i+1)D$ reaches the midpoint $(j+\tfrac12)d_x$.
Multiplying that comparison by two gives $2E+2D\geq d_x$. Our accumulator satisfies
$e=d_x-D-E$; substituting it gives $2e\leq d_x$, exactly the vertical decision.
The corresponding major/minor exchange gives the horizontal decision for steep
lines. The common two-test loop covers both orientations.

The word ``error'' names this decision state; it is not a distance measured in
pixels. Its useful invariant is
\[
 e=d_x-D-iD+jd_x,
 \qquad i=x-A_x,\quad j=s_y(y-A_y).
\]
A horizontal step subtracts $D$ from this expression, and a vertical step adds
$d_x$. The updates therefore keep it consistent with the coordinates. Replacing a
coordinate update without its corresponding accumulator update breaks that relation.

\section{Trace one saved decision through two updates}

For the default line, $d_x=5$, $d_y=-3$ and $s_y=1$, so the initial error is two.
Table~\ref{tab:revision-line-state} records state before each cell is painted.
The final cell is painted and then returned from; its listed error does not cause
another decision.

\begin{table}[htbp]
\centering
\begin{tabular}{rrrrll}
\toprule
$x$ & $y$ & $e$ & $2e$ & Coordinate steps & Next error\\
\midrule
1 & 1 & 2 & 4 & $x$ and $y$ & 4\\
2 & 2 & 4 & 8 & $x$ only & 1\\
3 & 2 & 1 & 2 & $x$ and $y$ & 3\\
4 & 3 & 3 & 6 & $x$ only & 0\\
5 & 3 & 0 & 0 & $x$ and $y$ & 2\\
6 & 4 & 2 & --- & return & ---\\
\bottomrule
\end{tabular}
\caption{A decision is frozen before either coordinate changes. The first update adds $-3$ to the error; a vertical update then adds $5$.}
\label{tab:revision-line-state}
\end{table}

The second row exposes a common mistake. Its saved doubled error is eight. The
horizontal update changes $e$ from four to one. If the vertical test recomputes
$2e$ at that point, it sees two instead of eight and incorrectly advances $y$.
The result would next paint $(3,3)$ rather than $(3,2)$. Saving the doubled error
is part of the algorithm, not a temporary-register convenience.

An \texttt{if}/\texttt{else} translation is also wrong: the first row requires both
updates. Think of the cursor as one state transition with two decisions derived
from the old state. An assembly program must preserve that relationship explicitly.

\section{Give the rasterizer a bounded interface}

The function \texttt{draw\_line} receives the address of four words in \texttt{\$a0}:
$x_0$, $y_0$, $x_1$, $y_1$, in that order. Each coordinate must be from zero through
seven. \texttt{\$a1} is the base address of 64 valid framebuffer words, and
\texttt{\$a2} contains the color, from zero through \texttt{0x00ffffff}. The function
returns the number of painted cells in \texttt{\$v0}. It preserves the input words,
unselected framebuffer words, saved registers and stack pointer. It may change
argument and temporary registers and \texttt{\$v1}.

Passing a pointer to the endpoint record keeps the four related coordinates
together without needing four separate argument registers in addition to the
frame and color. The function loads the record, then orders its working registers;
it never swaps the caller's memory. A source record and a mutable cursor are
different objects even when their values initially agree.

The complete program is \texttt{integer-line.asm} in \texttt{revision/labs/files}.
Its data and caller are below. Append the function listings that follow, in their
printed order. The caller saves the returned count before selecting the exit
service, since that service selector also uses \texttt{\$v0}.

\noindent\begin{minipage}{\linewidth}
\MARSCodeLabel{Supply a small frame and retain the result}
\begin{minted}{mars-idea}
# mars:tool bitmap width=64 height=64 unitWidth=8 unitHeight=8 base=static_data
.data
.align 2
frame:           .space 256
frame_guard:     .word 0x13579bdf
endpoints:       .word 1, 1, 6, 4
endpoints_guard: .word 0x2468ace0
color:           .word 0x00ffffff
color_guard:     .word 0x10203040
line_count:      .word 0
count_guard:     .word 0x31415926
.text
.globl main
main:
    dli $t0, -16
    and $sp, $sp, $t0
    la $a0, endpoints
    la $a1, frame
    lw $a2, color
    balc draw_line
line_complete:
    sw $v0, line_count
    li $v0, 10
    syscall
\end{minted}
\end{minipage}

\section{Translate the state without changing its meaning}

First load and order the endpoints. If the second $x$ is smaller, exchange both
coordinate pairs. If the $x$ values agree, use $y$ to decide the order. The two
swaps must keep each $x$ paired with its original $y$. \texttt{slt} writes one when
its first source is signed-less-than its second source, otherwise zero; the
following compact branch consumes that result. The \texttt{nop} instructions keep
the conditional branches' following instructions safe on fall-through paths, as
in Chapter~\ref{revision:decisions-and-repetition}.

\noindent\begin{minipage}{\linewidth}
\MARSCodeLabel{Order registers while preserving the endpoint record}
\begin{minted}{mars-idea}
draw_line:
    lw $t0, 0($a0)
    lw $t1, 4($a0)
    lw $t2, 8($a0)
    lw $t3, 12($a0)
    slt $t9, $t2, $t0
    bnezc $t9, swap_endpoints
    nop
    bnec $t0, $t2, ordered
    nop
    slt $t9, $t3, $t1
    beqzc $t9, ordered
    nop
swap_endpoints:
    move $v1, $t0
    move $t0, $t2
    move $t2, $v1
    move $v1, $t1
    move $t1, $t3
    move $t3, $v1
\end{minted}
\end{minipage}

After ordering, \texttt{\$t0}, \texttt{\$t1} hold the cursor and \texttt{\$t2},
\texttt{\$t3} the target. The next block puts $d_x$ in \texttt{\$t4}, nonpositive
$d_y$ in \texttt{\$t5}, $s_y$ in \texttt{\$t6} and the error in \texttt{\$t7}.
A falling line already has a negative raw vertical delta, so that path retains it
and changes the direction to $-1$. A nonnegative delta is negated. For a horizontal
line the direction is $+1$, but the decision rule never takes a vertical step.

\noindent\begin{minipage}{\linewidth}
\MARSCodeLabel{Initialize signed differences and the accumulator}
\begin{minted}{mars-idea}
ordered:
    subu $t4, $t2, $t0
    li $t6, 1
    subu $t5, $t3, $t1
    bgezc $t5, positive_delta
    nop
    li $t6, -1
    bc differences_ready
positive_delta:
    subu $t5, $zero, $t5
differences_ready:
    addu $t7, $t4, $t5
    move $v0, $zero
\end{minted}
\end{minipage}

The final block paints one cell at address $frame+4(8y+x)$. The first shift
multiplies $y$ by eight; the second multiplies the word index by four. The address
addition uses \texttt{daddu} to retain the full pointer width. Coordinate and error
arithmetic uses word instructions: every admitted value is small enough to be
represented exactly. For example, the invariant gives the conservative bound
$|e|\leq56$ for $0\leq i,j,d_x,D\leq7$, so doubling it cannot overflow a signed word.
This bound depends on the stated domain; it is not a license to accept arbitrary
32-bit endpoint differences.

\noindent\begin{minipage}{\linewidth}
\MARSCodeLabel{Paint, stop at the endpoint, then update from one snapshot}
\begin{minted}{mars-idea}
line_pixel:
    sll $v1, $t1, 3
    addu $v1, $v1, $t0
    dsll $v1, $v1, 2
    daddu $v1, $a1, $v1
    sw $a2, 0($v1)
    addiu $v0, $v0, 1
    bnec $t0, $t2, line_decide
    nop
    beqc $t1, $t3, line_return
    nop
line_decide:
    sll $t8, $t7, 1
    bltc $t8, $t5, skip_x
    nop
    addu $t7, $t7, $t5
    addiu $t0, $t0, 1
skip_x:
    bltc $t4, $t8, line_next
    nop
    addu $t7, $t7, $t4
    addu $t1, $t1, $t6
line_next:
    bc line_pixel
line_return:
    jrc $ra
\end{minted}
\end{minipage}

\texttt{bltc} compares signed register values and branches immediately when the
first is less than the second. The first test skips the horizontal update when
$e_2<d_y$. The second skips the vertical update when $d_x<e_2$. Both read
\texttt{\$t8}, which is written once per decision. The unconditional compact
\texttt{bc} returns to the next pixel calculation without a delay slot.

This routine is a leaf: it makes no further call, changes no saved register and
needs no stack frame. At \texttt{line\_return}, \texttt{\$ra} still holds its
caller's continuation. For the default input, \texttt{\$v0} is six. The finished
program leaves six in \texttt{line\_count}, exits normally and prints nothing.

\section{Check boundaries and adapt the operation}

Run the complete program with MIPS64 Release 6 and the Default layout. Open its
configured Bitmap Display. In Debug, pause at \texttt{line\_decide} on the second
visit: the current cell is $(2,2)$ and the current error is four. Step over the
snapshot and horizontal update, then verify that \texttt{\$t8} remains eight even
though \texttt{\$t7} becomes one. The pending vertical test must use eight.

Test a horizontal line, a vertical line, a steep falling line and coincident
endpoints. A point must paint once and return before any coordinate update. A
non-point line visits $\max(d_x,D)+1$ cells, at most eight in this domain. Its
monotone coordinate changes and nearest-cell decisions keep those cells inside
the endpoint bounding box, hence inside the framebuffer. At least one coordinate
advances per nonfinal iteration: failing both tests would require simultaneously
$e_2<d_y\leq0$ and $e_2>d_x\geq0$.

A halfway test checks a different issue from an endpoint test. Use $(0,0)$ to
$(2,1)$ and then reverse the inputs. Both images should contain $(1,1)$. Prefill
unselected cells with distinct nonzero words as well; a program that clears the
frame before drawing violates the overlay contract even if its white line looks
correct on an initially black display. Check the input record and guards, not just
the picture.

For independent work, draw alternating cells of the same canonical raster: paint
cell zero, preserve cell one, paint cell two, and continue. The sequence index
counts visits along the line, not the value of $x$, so vertical and steep lines
must work too. Return the number actually painted. The lab supplies a failing
starter and public checks. Keep the geometric decision state unchanged while
adding the state that controls whether a visited cell is stored.
